\section*{Capítulo \ref{ch:b2:elements-in-the-earth} --- Los elementos en la Tierra}

\begin{solution}{exo:b2:elements-in-the-earth:1}
El potasio es $2\times39.10/(2\times39.10+16.00) = 78.20/94.20 = 0.8301$ de la masa del
\ce{K2O}, de modo que hay $\num{5130}\times0.8301 = \qty{4259}{ppm}$ de K. El titanio es
$47.87/79.87 = 0.5993$ del \ce{TiO2}: $\num{27310}\times0.5993 = \qty{16370}{ppm}$ de Ti.
\end{solution}

\begin{solution}{exo:b2:elements-in-the-earth:2}
El aluminio es litófilo y refractario: la Tierra silicatada lo conservó todo, en proporciones casi
condríticas (1.16). El sodio es litófilo pero volátil (\qty{958}{K}): se
perdió antes de la acreción o durante ella, y quedó 0.22. El níquel condensa con el magnesio
(\qty{1353}{K}), de modo que la volatilidad no puede explicar 0.077: es siderófilo, y la mayor parte
se fue al núcleo. El iridio es refractario (\qty{1603}{K}) y muy siderófilo:
0.0029, casi todo él en el núcleo.
\end{solution}

\begin{solution}{exo:b2:elements-in-the-earth:3}
$\mathrm{La}_N = 31/0.232 = 133.6$, $\mathrm{Yb}_N = 2.0/0.163 = 12.27$, y
$(\mathrm{La}/\mathrm{Yb})_N = 10.9$. La corteza está enriquecida en todas las \omterm{def:b2:elements-in-the-earth:ree}{tierras raras}, y
once veces más en las ligeras que en las pesadas: está hecha de fundidos, que
toman con más facilidad las \omterm{def:b2:elements-in-the-earth:ree}{tierras raras} ligeras, incompatibles.
\end{solution}

\begin{solution}{exo:b2:elements-in-the-earth:4}
$\mathrm{Sm}_N = 4.7/0.145 = 32.4$, $\mathrm{Eu}_N = 1.0/0.0546 = 18.3$,
$\mathrm{Gd}_N = 4.0/0.198 = 20.2$; $\mathrm{Eu}^* = \sqrt{32.4\times20.2} = 25.6$ y
$\mathrm{Eu}/\mathrm{Eu}^* = 18.3/25.6 = 0.72$, una anomalía negativa. La plagioclasa, que
admite \ce{Eu^2+}, cristalizó a partir de los fundidos que construyeron la corteza superior y se quedó
a más profundidad.
\end{solution}

\begin{solution}{exo:b2:elements-in-the-earth:5}
$\bar D_{\mathrm{La}} = 0.65\times0.058+0.18\times0.001+0.17\times0.001 = 0.0377 +
0.0002 + 0.0002 = 0.038$; $\bar D_{\mathrm{Lu}} = 0.3445+1.044+0.0306 = 1.42$. El lutecio es
compatible ($\bar D > 1$), y el lantano muy incompatible. El granate contiene
$1.044/1.42 = 74$ por ciento del lutecio de la roca, aunque solo es el 18 por ciento de la
roca.
\end{solution}

\begin{solution}{exo:b2:elements-in-the-earth:6}
Para el sodio, la razón Na/Mg de la Tierra silicatada es $2670/\num{228000} = 0.01171$ y el valor
CI $5010/\num{95870} = 0.05226$: el cociente es 0.224. Para el níquel, $1960/\num{228000}
= 0.008596$ frente a $\num{10640}/\num{95870} = 0.1110$: 0.0775. El sodio está
empobrecido por su volatilidad (\qty{958}{K}); el níquel condensa por encima de \qty{1300}{K} y está
tres veces más empobrecido que el sodio porque es siderófilo: el núcleo se llevó la mayor parte
de él.
\end{solution}

\begin{solution}{exo:b2:elements-in-the-earth:7}
La: $5.21\times0.058 = \qty{0.302}{ppm}$; Lu: $0.53\times0.53 = \qty{0.281}{ppm}$. Los
valores de normalización se cancelan en el cociente de cocientes:
$(\mathrm{La}/\mathrm{Lu})_{N,\mathrm{cpx}}/(\mathrm{La}/\mathrm{Lu})_{N,\mathrm{melt}} =
D_{\mathrm{La}}/D_{\mathrm{Lu}} = 0.058/0.53 = 0.11$: el perfil del cristal es nueve
veces más inclinado hacia las \omterm{def:b2:elements-in-the-earth:ree}{tierras raras} pesadas que el del fundido.
\end{solution}

\begin{solution}{exo:b2:elements-in-the-earth:8}
Para la corteza continental superior, el cociente vale $(12/0.658)/(31/0.648) = 18.24/47.84 = 0.38$. Basalto hawaiano:
$(18.10/0.658)/(15.20/0.648) = 27.51/23.46 = 1.17$. MORB: $7.96/8.04 = 0.99$. Solo la
corteza, con un fuerte déficit de niobio, lleva la firma de las zonas de subducción.
\end{solution}

\begin{solution}{exo:b2:elements-in-the-earth:9}
\omterm{def:b2:elements-in-the-earth:bse}{Manto primitivo}: $240/0.0203 = \num{11800}$. MORB: K $= 0.160\times0.830\times10^4 =
\qty{1328}{ppm}$, K/U $= \num{11200}$. Corteza: K $= 2.80\times0.830\times10^4 =
\qty{23240}{ppm}$, K/U $= \num{8600}$. El potasio se multiplica por cien del manto a la corteza,
pero K/U cambia menos de un tercio: los dos elementos son más o menos igual de incompatibles,
y la fusión apenas los separa. El uranio es refractario, de modo que su abundancia se deduce
del modelo condrítico; K/U medido en lavas da entonces el potasio del planeta, que,
al ser volátil, no puede tomarse de las \omterm{def:b2:planet-formation-meteorites:chondrule}{condritas}.
\end{solution}

\begin{solution}{exo:b2:elements-in-the-earth:10}
En el equilibrio $C_{\mathrm{rock}} = \bar D\,C_{\mathrm{melt}}$ para cada elemento (y una
pequeña fracción de fundido deja la composición de la roca casi sin cambios), de modo que
$(C_1/C_2)_{\mathrm{melt}} = (C_{1}/\bar D_1)/(C_{2}/\bar D_2)$ para las $C_1,
C_2$ de la roca, y el cociente es $\bar D_2/\bar D_1$. Para La y Lu: $1.419/0.03805 = 37$. El fundido
de un manto con granate es 37 veces más rico en La respecto al Lu que su fuente: el granate
retuvo las \omterm{def:b2:elements-in-the-earth:ree}{tierras raras} pesadas, y el perfil del fundido es muy inclinado.
\end{solution}

\begin{solution}{exo:b2:elements-in-the-earth:11}
$C_{\mathrm{mix}} = fC_U + (1-f)C_M$ para cada elemento: lineal en $f$. El cociente
$\bigl(f\,\mathrm{Th}_U + (1-f)\mathrm{Th}_M\bigr)/\bigl(f\,\mathrm{Nb}_U +
(1-f)\mathrm{Nb}_M\bigr)$ es un cociente de dos funciones lineales, una hipérbola en $f$.
Th/Nb $= 0.2$ cuando $f(10.5-2.4) + (1-f)(0.404-1.048) = 0$, es decir $8.1f = 0.644(1-f)$,
$f = 0.644/8.744 = 0.074$: cerca de un 7 por ciento de corteza eleva Th/Nb de 0.077 a 0.2. Para
$f = 0.1$: Th $= 1.05 + 0.364 = 1.414$, Nb $= 1.2 + 4.716 = 5.916$, Th/Nb $= 0.24$.
\end{solution}

\begin{solution}{exo:b2:elements-in-the-earth:12}
Retirar un componente que no contiene ninguno de los dos elementos deja sin cambios ambas masas y, por tanto, su
cociente. Si se conservó todo el aluminio, $(\mathrm{Al}/\mathrm{Mg}) = 1.16$ veces CI
significa que la Tierra silicatada contiene $1/1.16 = 0.86$ de su magnesio condrítico: falta un 14 por
ciento. El silicio respecto al aluminio vale $0.83/1.16 = 0.72$ veces el valor CI:
faltaría un 28 por ciento del silicio condrítico.
\end{solution}

\begin{solution}{pb:b2:elements-in-the-earth:1}
\textbf{1.} Fe/Mg $= \num{182800}/\num{95870} = 1.907$; Ni/Mg $= 0.1110$; Si/Mg $=
1.111$; Al/Mg $= 0.0887$.

\textbf{2.} $(\num{182800}/55.85)/(\num{106500}/28.09) = 3273/3791 = 0.863$ átomos de hierro por
átomo de silicio, a menos de un 7 por ciento de los 0.81 de la fotosfera: razonable, ya que las \omterm{def:b2:elements-in-the-earth:ci}{condritas CI}
coinciden con el Sol en los elementos formadores de rocas.

\textbf{3.} La pérdida de un componente que no contiene ninguno de los dos elementos de una razón deja la
razón sin cambios (\cref{exo:b2:elements-in-the-earth:12}): solo cambian las concentraciones absolutas.

\textbf{4.} El magnesio es litófilo y no entra en el núcleo; es un elemento mayor,
bien determinado tanto en las \omterm{def:b2:planet-formation-meteorites:chondrule}{condritas} como en el manto. El hierro es el elemento que se reparte
entre núcleo y manto: no puede ser su propia referencia.

\textbf{5.} Los elementos que condensan juntos no fueron separados por la volatilidad en la
nebulosa, de modo que la Tierra global debería conservar sus razones. Al/Mg se aparta: el aluminio es
refractario (\qty{1653}{K}), y la razón Al/Mg de la Tierra silicatada vale 1.16 veces el valor CI.

\textbf{6.} Fe $= 8.05\times55.85/71.85 = 6.26$ por ciento; Ni $= 0.25\times58.69/74.69 =
0.196$ por ciento.

\textbf{7.} Fe/Mg $= 6.26/22.8 = 0.2746$; Ni/Mg $= 0.196/22.8 = 0.00860$.

\textbf{8.} Hierro: $1 - 0.2746/1.907 = 0.856$, falta un 85.6 por ciento; níquel:
$1 - 0.00860/0.1110 = 0.923$, un 92.3 por ciento.

\textbf{9.} El iridio es más siderófilo que el níquel, y el níquel más que el hierro:
cuanto más prefiere un elemento el metal, más se llevó de él el metal al núcleo.

\textbf{10.} En los silicatos, como \ce{Fe^2+} en el olivino y el piroxeno (8 por ciento de FeO):
el hierro oxidado a FeO se enlaza con el oxígeno como un \omterm{def:b2:elements-in-the-earth:goldschmidt}{elemento litófilo} y se quedó en el
manto; solo se hundió el hierro metálico.

\textbf{11.} $(1-x)C_{\mathrm{Fe}} + xc_{\mathrm{Fe}} = r_{\mathrm{Fe}}(1-x)C_{\mathrm{Mg}}$;
dividiendo por $(1-x)c_{\mathrm{Fe}}$ se obtiene la proposición.

\textbf{12.} $r_{\mathrm{Fe}}C_{\mathrm{Mg}} - C_{\mathrm{Fe}} = 1.907\times22.8 - 6.26 =
37.22$; $y = 37.22/85.5 = 0.4353$, $x = y/(1+y) = 0.303$.

\textbf{13.} $0.1110\times22.8 - 0.196 = 2.335$; $y = 2.335/5.2 = 0.4490$, $x = 0.310$.

\textbf{14.} $y = (85.5\times37.22 + 5.2\times2.335)/(85.5^2 + 5.2^2) =
3194.4/7337.3 = 0.4354$, $x = 0.303$.

\textbf{15.} Cada ecuación pesa como el cuadrado de su coeficiente: $85.5^2/5.2^2 = 270$
veces más para el hierro, el elemento mucho más abundante.

\textbf{16.} $37.22/2.335 = 15.9$, frente a $85.5/5.2 = 16.4$ en el modelo: una diferencia del 3 por
ciento.

\textbf{17.} $0.303\times\num{5.972e24} = \qty{1.81e24}{kg}$.

\textbf{18.} $(32.5-30.3)/32.5 = 7$ por ciento menos.

\textbf{19.} $y = 0.325/0.675 = 0.4815$, Fe en el núcleo $= 37.22/0.4815 = 77.3$ por ciento,
y Ni $= 2.335/0.4815 = 4.85$ por ciento: un 18 por ciento de elementos ligeros, cerca del doble del 5 al
10 por ciento que permite la densidad. La razón Fe/Mg de la Tierra debe ser algo mayor que el valor
CI, o su magnesio silicatado mayor que el del modelo.

\textbf{20.} Silicio, azufre, carbono, oxígeno e hidrógeno. Al silicio le falta en la
Tierra silicatada: su razón normalizada al Mg, 0.83, está por debajo del nivel refractario.

\textbf{21.} $1.111\times22.8 - 21.0 = 4.331$; $4.331/0.4354 = 9.9$ por ciento de Si en el
núcleo.

\textbf{22.} 9.9 frente a 6 por ciento: el núcleo puede contener más silicio que el modelo, o
parte del silicio nunca se acretó. El silicio condensa a \qty{1310}{K}, un poco
por debajo del magnesio y del hierro (1336 y \qty{1334}{K}): parte pudo quedarse en el gas.

\textbf{23.} Núcleo: $0.325\times0.5 = \qty{0.1625}{ppm}$ de la masa de la Tierra; Tierra
silicatada: $0.675\times0.001 = \qty{0.000675}{ppm}$. El núcleo contiene
$0.1625/0.1632 = 99.6$ por ciento del oro.

\textbf{24.} La química sola sitúa
\textbf{\qty[text-series-to-math=true]{30.3}{\percent}} de la masa de la Tierra en su núcleo,
a menos de un 7 por ciento del 32.5 por ciento sismológico.
\end{solution}
